\section{Electromagnetism}
\label{sec:electromagnetism}

\noindent\emph{Dense reference extracted from Tong's \texttt{tong/electro.pdf}.}
SI units throughout unless stated. Fields live on $\R^3\times\R$; repeated Cartesian
or spacetime indices are summed where clear. Relativistic convention:
$X^\mu=(ct,\mathbf x)$ and $\eta_{\mu\nu}=\operatorname{diag}(+,-,-,-)$.

\subsection{Introduction and Maxwell Data}

\subsubsection*{Charge, current, and conservation}
\[
q\in\R,\qquad [q]=\mathrm C,\qquad q=ne,\quad e=1.602176634\times10^{-19}\,\mathrm C.
\]
\[
Q_V(t)=\int_V\rho(\mathbf x,t)\,\dd^3x,\qquad
I_S=\int_S\mathbf J\cdot\dd\mathbf S,\qquad
\rho_q(\mathbf x,t)=q\,\delta^{(3)}(\mathbf x-\mathbf r(t)).
\]
\[
\boxed{\frac{\pd\rho}{\pd t}+\nabla\cdot\mathbf J=0},\qquad
\frac{\dd Q_V}{\dd t}=-\int_{\pd V}\mathbf J\cdot\dd\mathbf S.
\]
Continuity is the compatibility condition for admissible sources $(\rho,\mathbf J)$.
\src{Tong EM \S 1.1}

\subsubsection*{Lorentz force and Maxwell equations}
\[
\boxed{\mathbf F=q(\mathbf E+\mathbf v\times\mathbf B)},\qquad
\mathbf f=\rho\mathbf E+\mathbf J\times\mathbf B.
\]
\[
\boxed{
\begin{aligned}
\nabla\cdot\mathbf E&=\frac{\rho}{\epsilon_0},&
\nabla\cdot\mathbf B&=0,\\
\nabla\times\mathbf E&=-\frac{\pd\mathbf B}{\pd t},&
\nabla\times\mathbf B&=\mu_0\mathbf J+\mu_0\epsilon_0\frac{\pd\mathbf E}{\pd t}.
\end{aligned}}
\]
\[
c=\frac{1}{\sqrt{\epsilon_0\mu_0}},\qquad
\nabla\cdot(\nabla\times\mathbf B)=0
\Longrightarrow
\frac{\pd\rho}{\pd t}+\nabla\cdot\mathbf J=0.
\]
The displacement current term is forced by charge conservation. \src{Tong EM \S 1.2}

\subsection{Electrostatics}

\subsubsection*{Gauss law and Coulomb templates}
Static regime:
\[
\pd_t\rho=0,\qquad \mathbf J=\mathbf0,\qquad
\boxed{\nabla\cdot\mathbf E=\rho/\epsilon_0},\qquad
\boxed{\nabla\times\mathbf E=\mathbf0}.
\]
\[
\oint_{\pd V}\mathbf E\cdot\dd\mathbf S=\frac{Q_{\rm enc}}{\epsilon_0}.
\]
Point charge and force:
\[
\boxed{\mathbf E(\mathbf r)=\frac{q}{4\pi\epsilon_0r^2}\hat{\mathbf r}},\qquad
\boxed{\mathbf F_{12}=\frac{q_1q_2}{4\pi\epsilon_0|\mathbf r_1-\mathbf r_2|^2}
\widehat{(\mathbf r_1-\mathbf r_2)}}.
\]
Spherical charge of total $Q$ inside $R$:
\[
\mathbf E(r>R)=\frac{Q}{4\pi\epsilon_0r^2}\hat{\mathbf r}.
\]
Uniform solid sphere:
\[
\mathbf E(r)=
\begin{cases}
\displaystyle \frac{Qr}{4\pi\epsilon_0R^3}\hat{\mathbf r}
=\frac{\rho r}{3\epsilon_0}\hat{\mathbf r},&r<R,\\[0.6em]
\displaystyle \frac{Q}{4\pi\epsilon_0r^2}\hat{\mathbf r},&r>R.
\end{cases}
\]
Infinite line charge $\eta$ and plane charge $\sigma$:
\[
\mathbf E_{\rm line}(s)=\frac{\eta}{2\pi\epsilon_0s}\hat{\mathbf s},\qquad
\mathbf E_{\rm plane}(z>0)=\frac{\sigma}{2\epsilon_0}\hat{\mathbf z},\quad
\mathbf E_{\rm plane}(z<0)=-\frac{\sigma}{2\epsilon_0}\hat{\mathbf z}.
\]
\src{Tong EM \S 2.1}

\subsubsection*{Surface charge discontinuities}
For unit normal $\hat{\mathbf n}$ from side $-$ to $+$:
\[
\boxed{\hat{\mathbf n}\cdot(\mathbf E_+-\mathbf E_-)=\sigma/\epsilon_0},\qquad
\boxed{\hat{\mathbf n}\times(\mathbf E_+-\mathbf E_-)=\mathbf0}.
\]
Equivalently, if $\mathbf E=-\nabla\phi$,
\[
\phi_+=\phi_-,\qquad
\pd_n\phi_+-\pd_n\phi_-=-\sigma/\epsilon_0.
\]
\src{Tong EM \S 2.1.4}

\subsubsection*{Potential, Green function, and multipoles}
\[
\mathbf E=-\nabla\phi,\qquad
\boxed{\nabla^2\phi=-\rho/\epsilon_0}.
\]
With $\phi(\infty)=0$ for compact sources:
\[
\boxed{\phi(\mathbf r)=\frac{1}{4\pi\epsilon_0}\int
\frac{\rho(\mathbf r')}{|\mathbf r-\mathbf r'|}\,\dd^3r'},\qquad
\nabla^2\frac{1}{|\mathbf r-\mathbf r'|}=-4\pi\delta^{(3)}(\mathbf r-\mathbf r').
\]
Point charge and dipole:
\[
\phi_q=\frac{q}{4\pi\epsilon_0r},\qquad
\mathbf p=\sum_i q_i\mathbf r_i,\qquad
\boxed{\phi_{\rm dip}(\mathbf r)=\frac{\mathbf p\cdot\hat{\mathbf r}}{4\pi\epsilon_0r^2}},
\]
\[
\boxed{\mathbf E_{\rm dip}(\mathbf r)=\frac{1}{4\pi\epsilon_0r^3}
\left[3(\mathbf p\cdot\hat{\mathbf r})\hat{\mathbf r}-\mathbf p\right]}.
\]
Large-$r$ expansion:
\[
\phi(\mathbf r)=\frac{1}{4\pi\epsilon_0}\left[
\frac{Q}{r}+\frac{\mathbf p\cdot\hat{\mathbf r}}{r^2}
\frac{1}{2r^3}\sum_{ij}Q_{ij}\hat r_i\hat r_j+\cdots\right],
\]
\[
Q=\int\rho\,\dd^3r,\qquad
\mathbf p=\int\rho\,\mathbf r\,\dd^3r,\qquad
Q_{ij}=\int\rho(3x_ix_j-r^2\delta_{ij})\,\dd^3r.
\]
If $Q=0$, $\mathbf p$ is origin independent. \src{Tong EM \S 2.2}

\subsubsection*{Energy and forces}
Assembly energy:
\[
U=\frac12\int\rho\phi\,\dd^3x
=\frac{\epsilon_0}{2}\int |\mathbf E|^2\,\dd^3x
\quad\text{for fields decaying at infinity}.
\]
Discrete charges:
\[
U=\frac{1}{4\pi\epsilon_0}\sum_{i<j}\frac{q_iq_j}{|\mathbf r_i-\mathbf r_j|},
\]
excluding divergent point self-energies. Dipole in external field:
\[
U=-\mathbf p\cdot\mathbf E,\qquad
\boldsymbol\tau=\mathbf p\times\mathbf E,\qquad
\mathbf F=\nabla(\mathbf p\cdot\mathbf E)\quad(\mathbf p\ \text{fixed}).
\]
Dipole-dipole energy:
\[
U_{12}=\frac{1}{4\pi\epsilon_0r^3}
\left[\mathbf p_1\cdot\mathbf p_2-3(\mathbf p_1\cdot\hat{\mathbf r})
(\mathbf p_2\cdot\hat{\mathbf r})\right].
\]
\src{Tong EM \S 2.3}

\subsubsection*{Conductors, capacitance, and boundary value problems}
Electrostatic conductor:
\[
\mathbf E_{\rm inside}=0,\qquad \rho_{\rm inside}=0,\qquad
\phi=\text{constant on each conductor},\qquad
\mathbf E_{\rm surface}=\frac{\sigma}{\epsilon_0}\hat{\mathbf n}.
\]
Capacitance:
\[
Q=CV,\qquad U=\frac12 CV^2=\frac{Q^2}{2C}.
\]
Parallel plate, sphere, coaxial cylinders, concentric spheres:
\[
C_{\rm plate}=\frac{\epsilon_0A}{d},\qquad
C_{\rm sphere}=4\pi\epsilon_0R,\qquad
\frac{C_{\rm coax}}{L}=\frac{2\pi\epsilon_0}{\log(b/a)},\qquad
C_{\rm sph}=\frac{4\pi\epsilon_0ab}{b-a}.
\]
Uniqueness: for $\nabla^2\phi=-\rho/\epsilon_0$ on $\Omega$, Dirichlet data on
$\pd\Omega$ determine $\phi$ uniquely; Neumann data determine $\phi$ up to a constant
subject to the total-charge compatibility condition. Proof kernel:
\[
f=\phi_1-\phi_2,\quad \nabla^2f=0,\quad
\int_\Omega|\nabla f|^2\,\dd^3x=\int_{\pd\Omega}f\,\pd_nf\,\dd S.
\]
Method of images replaces conductors by fictitious charges outside the physical region
while preserving boundary data. Grounded plane, charge $q$ at height $a$:
\[
q_{\rm image}=-q\ \text{at }z=-a,\qquad
F_z=-\frac{q^2}{16\pi\epsilon_0a^2}.
\]
Grounded sphere radius $R$, external charge $q$ at distance $a>R$:
\[
q'=-\frac{R}{a}q,\qquad b=\frac{R^2}{a}.
\]
\src{Tong EM \S 2.4}

\subsection{Magnetostatics}

\subsubsection*{Ampere law and surface currents}
Static neutral currents:
\[
\nabla\cdot\mathbf J=0,\qquad
\boxed{\nabla\times\mathbf B=\mu_0\mathbf J},\qquad
\boxed{\nabla\cdot\mathbf B=0}.
\]
\[
\oint_C\mathbf B\cdot\dd\mathbf r=\mu_0I_{\rm enc}.
\]
Straight wire, sheet current, solenoid:
\[
\mathbf B_{\rm wire}(s)=\frac{\mu_0I}{2\pi s}\hat{\boldsymbol\phi},
\]
\[
\mathbf K=K\hat{\mathbf x}\ \text{on }z=0:\quad
\mathbf B(z>0)=-\frac{\mu_0K}{2}\hat{\mathbf y},\quad
\mathbf B(z<0)=+\frac{\mu_0K}{2}\hat{\mathbf y},
\]
\[
\text{long solenoid:}\qquad \mathbf B_{\rm in}=\mu_0nI\,\hat{\mathbf z},\qquad
\mathbf B_{\rm out}\simeq0.
\]
Boundary conditions at surface current $\mathbf K$:
\[
\boxed{\hat{\mathbf n}\cdot(\mathbf B_+-\mathbf B_-)=0},\qquad
\boxed{\hat{\mathbf n}\times(\mathbf B_+-\mathbf B_-)=\mu_0\mathbf K}.
\]
\src{Tong EM \S 3.1}

\subsubsection*{Vector potential, gauge, Biot-Savart}
\[
\boxed{\mathbf B=\nabla\times\mathbf A},\qquad
\mathbf A\mapsto\mathbf A+\nabla\chi.
\]
Coulomb gauge $\nabla\cdot\mathbf A=0$:
\[
\boxed{\nabla^2\mathbf A=-\mu_0\mathbf J},\qquad
\boxed{\mathbf A(\mathbf r)=\frac{\mu_0}{4\pi}
\int\frac{\mathbf J(\mathbf r')}{|\mathbf r-\mathbf r'|}\,\dd^3r'}.
\]
Biot-Savart:
\[
\boxed{\mathbf B(\mathbf r)=\frac{\mu_0}{4\pi}\int
\frac{\mathbf J(\mathbf r')\times(\mathbf r-\mathbf r')}
{|\mathbf r-\mathbf r'|^3}\,\dd^3r'}.
\]
Thin wire:
\[
\mathbf B(\mathbf r)=\frac{\mu_0I}{4\pi}\int_C
\frac{\dd\mathbf r'\times(\mathbf r-\mathbf r')}{|\mathbf r-\mathbf r'|^3}.
\]
\src{Tong EM \S 3.2.3}

\subsubsection*{Monopoles and linking number}
If magnetic charge $g$ existed,
\[
\nabla\cdot\mathbf B=\mu_0\rho_m,\qquad
\mathbf B_g=\frac{\mu_0g}{4\pi r^2}\hat{\mathbf r},
\]
so $\mathbf A$ cannot be globally smooth. Dirac quantisation:
\[
\boxed{\frac{qg}{2\pi\hbar}\in\mathbb Z}.
\]
For two closed curves $C,C'$:
\[
\boxed{L(C,C')=\frac{1}{4\pi}\oint_C\oint_{C'}
\frac{(\dd\mathbf r\times\dd\mathbf r')\cdot(\mathbf r-\mathbf r')}
{|\mathbf r-\mathbf r'|^3}\in\mathbb Z}.
\]
\src{Tong EM \S 3.2.1--3.2.4}

\subsubsection*{Magnetic dipoles and forces}
Current loop:
\[
\mathbf m=IA\,\hat{\mathbf n},\qquad
\mathbf A(\mathbf r)=\frac{\mu_0}{4\pi}\frac{\mathbf m\times\mathbf r}{r^3},\qquad
\mathbf B(\mathbf r)=\frac{\mu_0}{4\pi r^3}\left[3(\mathbf m\cdot\hat{\mathbf r})\hat{\mathbf r}-\mathbf m\right].
\]
General localized current:
\[
\mathbf m=\frac12\int\mathbf r\times\mathbf J(\mathbf r)\,\dd^3r.
\]
Force on current distribution:
\[
\mathbf F=\int\mathbf J\times\mathbf B\,\dd^3x,\qquad
\dd\mathbf F=I\,\dd\mathbf r\times\mathbf B.
\]
Between two closed wires:
\[
\mathbf F_{12}=-\frac{\mu_0I_1I_2}{4\pi}\oint_{C_1}\oint_{C_2}
\frac{(\dd\mathbf r_1\cdot\dd\mathbf r_2)(\mathbf r_1-\mathbf r_2)}
{|\mathbf r_1-\mathbf r_2|^3}.
\]
Dipole in external field:
\[
U=-\mathbf m\cdot\mathbf B,\qquad
\boldsymbol\tau=\mathbf m\times\mathbf B,\qquad
\mathbf F=\nabla(\mathbf m\cdot\mathbf B).
\]
Electron scale:
\[
\boldsymbol\mu=-g\frac{e}{2m_e}\mathbf S,\qquad g\simeq2.
\]
\src{Tong EM \S 3.3--3.4}

\subsubsection*{Unit systems}
\[
\alpha=\frac{e^2}{4\pi\epsilon_0\hbar c}\approx\frac{1}{137}.
\]
SI keeps $\epsilon_0,\mu_0$ explicit; Lorentz-Heaviside/natural conventions absorb
these constants and often set $c=\hbar=1$. \src{Tong EM \S 3.5}

\subsection{Electrodynamics}

\subsubsection*{Faraday induction and moving circuits}
\[
\Phi_B=\int_S\mathbf B\cdot\dd\mathbf S,\qquad
\mathcal E=\oint_C\mathbf E\cdot\dd\mathbf r,\qquad
\boxed{\mathcal E=-\frac{\dd\Phi_B}{\dd t}}\quad(C,S\ \text{fixed}).
\]
Moving wire:
\[
\boxed{\mathcal E=\oint_{C(t)}(\mathbf E+\mathbf v\times\mathbf B)\cdot\dd\mathbf r
=-\frac{\dd}{\dd t}\int_{S(t)}\mathbf B\cdot\dd\mathbf S}.
\]
Lenz sign: induced current opposes the flux change. \src{Tong EM \S 4.1}

\subsubsection*{Inductance, field energy, and resistance}
For linear circuits:
\[
\Phi_i=\sum_jL_{ij}I_j,\qquad L_{ij}=L_{ji},\qquad
\mathcal E_i=-\sum_jL_{ij}\dot I_j.
\]
Single inductor:
\[
\Phi=LI,\qquad U_L=\frac12LI^2.
\]
Magnetic field energy:
\[
\boxed{U_B=\frac{1}{2\mu_0}\int|\mathbf B|^2\,\dd^3x}.
\]
Ohm law and Joule heating:
\[
\mathbf J=\sigma\mathbf E,\qquad \rho_{\rm res}=1/\sigma,\qquad
V=IR,\qquad R=\rho_{\rm res}\frac{\ell}{A},
\]
\[
P=IV=I^2R=\frac{V^2}{R},\qquad
\frac{\dd U_{\rm heat}}{\dd t}=\int\mathbf J\cdot\mathbf E\,\dd^3x.
\]
\src{Tong EM \S 4.1.2--4.1.3}

\subsubsection*{Displacement current and waves}
Ampere correction:
\[
\boxed{\nabla\times\mathbf B=\mu_0\mathbf J+\mu_0\epsilon_0\frac{\pd\mathbf E}{\pd t}}.
\]
Vacuum wave equations:
\[
\boxed{\left(\frac{1}{c^2}\pd_t^2-\nabla^2\right)\mathbf E=0},\qquad
\boxed{\left(\frac{1}{c^2}\pd_t^2-\nabla^2\right)\mathbf B=0}.
\]
Plane wave:
\[
\mathbf E=\mathbf E_0e^{i(\mathbf k\cdot\mathbf x-\omega t)},\qquad
\mathbf B=\mathbf B_0e^{i(\mathbf k\cdot\mathbf x-\omega t)},\qquad
\omega=c|\mathbf k|,
\]
\[
\mathbf k\cdot\mathbf E_0=\mathbf k\cdot\mathbf B_0=0,\qquad
\boxed{\mathbf B_0=\frac{1}{c}\hat{\mathbf k}\times\mathbf E_0}.
\]
Linear and circular polarisations:
\[
\mathbf E_0=E_1\mathbf e_1+E_2e^{i\delta}\mathbf e_2,\qquad
\delta=0,\pi:\ \text{linear},\quad \delta=\pm\pi/2,\ |E_1|=|E_2|:\ \text{circular}.
\]
\src{Tong EM \S 4.2--4.3}

\subsubsection*{Reflection from a perfect conductor}
At a perfect conductor surface:
\[
\mathbf E_{\parallel}=0,\qquad \mathbf B_\perp=0.
\]
Normal incidence on $z=0$, conductor $z>0$:
\[
\mathbf E_I=E_0\hat{\mathbf x}e^{i(kz-\omega t)},\qquad
\mathbf E_R=-E_0\hat{\mathbf x}e^{i(-kz-\omega t)}.
\]
Outside field:
\[
\mathbf E=2iE_0\hat{\mathbf x}\,e^{-i\omega t}\sin kz,\qquad
\mathbf B=\frac{2E_0}{c}\hat{\mathbf y}\,e^{-i\omega t}\cos kz.
\]
Surface current supports the tangential magnetic discontinuity:
\[
\hat{\mathbf n}\times(\mathbf B_{\rm out}-\mathbf B_{\rm in})=\mu_0\mathbf K.
\]
\src{Tong EM \S 4.3.3}

\subsubsection*{Poynting theorem}
\[
u_{\rm em}=\frac{\epsilon_0}{2}|\mathbf E|^2+\frac{1}{2\mu_0}|\mathbf B|^2,\qquad
\boxed{\mathbf S=\frac{1}{\mu_0}\mathbf E\times\mathbf B}.
\]
\[
\boxed{\frac{\pd u_{\rm em}}{\pd t}+\nabla\cdot\mathbf S=-\mathbf J\cdot\mathbf E}.
\]
For matter energy obeying $\pd_tu_{\rm matter}=\mathbf J\cdot\mathbf E$:
\[
\pd_t(u_{\rm em}+u_{\rm matter})+\nabla\cdot\mathbf S=0.
\]
\src{Tong EM \S 4.4}

\subsection{Electromagnetism and Relativity}

\subsubsection*{Four-vectors and current}
\[
X^\mu=(ct,\mathbf x),\qquad
\pd_\mu=\left(\frac{1}{c}\pd_t,\nabla\right),\qquad
\pd^\mu=\left(\frac{1}{c}\pd_t,-\nabla\right).
\]
Proper time and four-velocity:
\[
\dd\tau=\frac{\dd t}{\gamma},\qquad
\gamma=(1-v^2/c^2)^{-1/2},\qquad
U^\mu=\frac{\dd X^\mu}{\dd\tau}=\gamma(c,\mathbf v).
\]
Four-current:
\[
\boxed{J^\mu=(c\rho,\mathbf J)},\qquad
\boxed{\pd_\mu J^\mu=0}.
\]
\src{Tong EM \S 5.1--5.2}

\subsubsection*{Gauge potential and field tensor}
\[
A^\mu=(\phi/c,\mathbf A),\qquad
\mathbf E=-\nabla\phi-\pd_t\mathbf A,\qquad
\mathbf B=\nabla\times\mathbf A.
\]
Gauge symmetry:
\[
\mathbf A\mapsto\mathbf A+\nabla\chi,\qquad
\phi\mapsto\phi-\pd_t\chi,\qquad
A_\mu\mapsto A_\mu-\pd_\mu\chi.
\]
\[
\boxed{F_{\mu\nu}=\pd_\mu A_\nu-\pd_\nu A_\mu},\qquad
F^{0i}=-E_i/c,\qquad F^{ij}=-\epsilon^{ijk}B_k
\]
with sign conventions fixed by the metric above. Dual:
\[
\widetilde F^{\mu\nu}=\frac12\epsilon^{\mu\nu\rho\sigma}F_{\rho\sigma}.
\]
Lorentz scalars:
\[
\boxed{F_{\mu\nu}F^{\mu\nu}=2(B^2-E^2/c^2)},\qquad
\boxed{\widetilde F_{\mu\nu}F^{\mu\nu}=-4\mathbf E\cdot\mathbf B/c}.
\]
Boost along $x$:
\[
E_\parallel'=E_\parallel,\quad B_\parallel'=B_\parallel,\qquad
\mathbf E'_\perp=\gamma(\mathbf E+\mathbf v\times\mathbf B)_\perp,
\]
\[
\mathbf B'_\perp=\gamma\left(\mathbf B-\frac{\mathbf v\times\mathbf E}{c^2}\right)_\perp .
\]
\src{Tong EM \S 5.3}

\subsubsection*{Covariant Maxwell equations and Lorentz force}
\[
\boxed{\pd_\mu F^{\mu\nu}=\mu_0J^\nu},\qquad
\boxed{\pd_\mu\widetilde F^{\mu\nu}=0}.
\]
Equivalently:
\[
\pd_\rho F_{\mu\nu}+\pd_\nu F_{\rho\mu}+\pd_\mu F_{\nu\rho}=0.
\]
Particle equation:
\[
p^\mu=mU^\mu,\qquad
\boxed{\frac{\dd p^\mu}{\dd\tau}=qF^{\mu\nu}U_\nu}.
\]
Spatial part gives $\dd(\gamma m\mathbf v)/\dd t=q(\mathbf E+\mathbf v\times\mathbf B)$.
\src{Tong EM \S 5.4}

\subsubsection*{Actions and stress tensor}
Nonrelativistic charged particle:
\[
S=\int\left(\frac12m\dot{\mathbf x}^2-q\phi+q\dot{\mathbf x}\cdot\mathbf A\right)\dd t,
\quad
\mathbf p_{\rm canonical}=m\dot{\mathbf x}+q\mathbf A.
\]
Relativistic particle plus interaction:
\[
S_{\rm p}=-mc^2\int\dd\tau,\qquad
S_{\rm int}=-\int J^\mu A_\mu\,\dd^4x.
\]
Maxwell action:
\[
\boxed{S_{\rm EM}=-\frac{1}{4\mu_0}\int F_{\mu\nu}F^{\mu\nu}\,\dd^4x
-\int J^\mu A_\mu\,\dd^4x}
\]
gives $\pd_\mu F^{\mu\nu}=\mu_0J^\nu$. Stress tensor:
\[
\boxed{T^{\mu\nu}=\frac{1}{\mu_0}
\left(-F^{\mu\rho}F^\nu{}_\rho+\frac14\eta^{\mu\nu}F_{\rho\sigma}F^{\rho\sigma}\right)}.
\]
Components:
\[
T^{00}=u_{\rm em},\qquad T^{0i}=c\,g_i,\qquad
\mathbf g=\epsilon_0\mathbf E\times\mathbf B=\mathbf S/c^2,
\]
\[
\sigma_{ij}=\epsilon_0\left(E_iE_j-\frac12\delta_{ij}E^2\right)
\frac{1}{\mu_0}\left(B_iB_j-\frac12\delta_{ij}B^2\right).
\]
Local balance with sources:
\[
\pd_\mu T^{\mu\nu}=-F^{\nu\rho}J_\rho.
\]
\src{Tong EM \S 5.5--5.6}

\subsection{Electromagnetic Radiation}

\subsubsection*{Retarded potentials and Green functions}
Lorenz gauge:
\[
\boxed{\pd_\mu A^\mu=0}\quad\Longleftrightarrow\quad
\nabla\cdot\mathbf A+\frac{1}{c^2}\pd_t\phi=0.
\]
Wave equations:
\[
\left(\nabla^2-\frac{1}{c^2}\pd_t^2\right)\phi=-\frac{\rho}{\epsilon_0},\qquad
\left(\nabla^2-\frac{1}{c^2}\pd_t^2\right)\mathbf A=-\mu_0\mathbf J.
\]
Retarded Green function:
\[
G_{\rm ret}(\mathbf x,t;\mathbf x',t')
=\frac{\delta(t'-t+|\mathbf x-\mathbf x'|/c)}
{4\pi|\mathbf x-\mathbf x'|}.
\]
Retarded potentials:
\[
\boxed{\phi(\mathbf x,t)=\frac{1}{4\pi\epsilon_0}\int
\frac{\rho(\mathbf x',t_r)}{|\mathbf x-\mathbf x'|}\,\dd^3x'},\qquad
\boxed{\mathbf A(\mathbf x,t)=\frac{\mu_0}{4\pi}\int
\frac{\mathbf J(\mathbf x',t_r)}{|\mathbf x-\mathbf x'|}\,\dd^3x'},
\]
\[
t_r=t-\frac{|\mathbf x-\mathbf x'|}{c}.
\]
\src{Tong EM \S 6.1}

\subsubsection*{Dipole radiation}
Localized source size $d\ll\lambda$, observer $r\gg d$, $\hat{\mathbf n}=\mathbf x/r$:
\[
\mathbf p(t)=\int\rho(\mathbf x,t)\mathbf x\,\dd^3x.
\]
Radiation-zone fields:
\[
\boxed{\mathbf E_{\rm rad}(\mathbf x,t)=\frac{1}{4\pi\epsilon_0c^2r}
\hat{\mathbf n}\times\big[\hat{\mathbf n}\times\ddot{\mathbf p}(t-r/c)\big]},
\]
\[
\boxed{\mathbf B_{\rm rad}=\frac{1}{c}\hat{\mathbf n}\times\mathbf E_{\rm rad}}.
\]
Angular and total power:
\[
\frac{\dd P}{\dd\Omega}=
\frac{|\ddot{\mathbf p}|^2\sin^2\theta}{16\pi^2\epsilon_0c^3},\qquad
\boxed{P=\frac{|\ddot{\mathbf p}|^2}{6\pi\epsilon_0c^3}}.
\]
For $\mathbf p(t)=\mathbf p_0\cos\omega t$:
\[
\langle P\rangle=\frac{\omega^4p_0^2}{12\pi\epsilon_0c^3}.
\]
Nonrelativistic point charge, $\mathbf p=q\mathbf r$:
\[
\boxed{P_{\rm Larmor}=\frac{q^2a^2}{6\pi\epsilon_0c^3}}.
\]
Next multipoles:
\[
P_{\rm mag\ dip}=\frac{\mu_0|\ddot{\mathbf m}|^2}{6\pi c^3},\qquad
P_{\rm elec\ quad}\sim \frac{1}{c^5}\left(\dddot Q_{ij}\right)^2.
\]
\src{Tong EM \S 6.2}

\subsubsection*{Scattering}
Driven nonrelativistic charge:
\[
m\ddot{\mathbf x}=q\mathbf E_0e^{-i\omega t},\qquad
\ddot{\mathbf p}=q\ddot{\mathbf x}=\frac{q^2}{m}\mathbf E_0e^{-i\omega t}.
\]
Thomson differential cross-section for incident polarisation $\boldsymbol\epsilon$:
\[
\boxed{\frac{\dd\sigma}{\dd\Omega}=
\left(\frac{q^2}{4\pi\epsilon_0mc^2}\right)^2
|\hat{\mathbf n}\times(\hat{\mathbf n}\times\boldsymbol\epsilon)|^2}.
\]
Unpolarised:
\[
\frac{\dd\sigma}{\dd\Omega}=\frac{r_e^2}{2}(1+\cos^2\theta),\qquad
\sigma_T=\frac{8\pi}{3}r_e^2,\qquad
r_e=\frac{e^2}{4\pi\epsilon_0m_ec^2}.
\]
Rayleigh scattering for bound charge with natural frequency $\omega_0$:
\[
\sigma_{\rm Rayleigh}\propto\omega^4\quad(\omega\ll\omega_0).
\]
\src{Tong EM \S 6.3}

\subsubsection*{Lienard-Wiechert fields}
Point charge worldline $\mathbf r(t)$; retarded time $t_\ast$:
\[
t_\ast+\frac{|\mathbf x-\mathbf r(t_\ast)|}{c}=t,\qquad
\mathbf R=\mathbf x-\mathbf r(t_\ast),\quad
\mathbf n=\mathbf R/R,\quad \boldsymbol\beta=\mathbf v(t_\ast)/c.
\]
Potentials:
\[
\boxed{\phi=\frac{q}{4\pi\epsilon_0}\frac{1}{R(1-\mathbf n\cdot\boldsymbol\beta)}},\qquad
\boxed{\mathbf A=\frac{\mu_0q}{4\pi}\frac{\mathbf v}{R(1-\mathbf n\cdot\boldsymbol\beta)}}.
\]
Fields:
\[
\mathbf E=\frac{q}{4\pi\epsilon_0}\left[
\frac{(1-\beta^2)(\mathbf n-\boldsymbol\beta)}
{R^2(1-\mathbf n\cdot\boldsymbol\beta)^3}
\frac{\mathbf n\times\{(\mathbf n-\boldsymbol\beta)\times\dot{\boldsymbol\beta}\}}
{cR(1-\mathbf n\cdot\boldsymbol\beta)^3}\right]_{t_\ast},
\]
\[
\boxed{\mathbf B=\frac{1}{c}\mathbf n\times\mathbf E}.
\]
First term is velocity/Coulomb field $\sim R^{-2}$; second is radiation field $\sim R^{-1}$.
Relativistic radiated power:
\[
\boxed{P=\frac{q^2\gamma^6}{6\pi\epsilon_0c^3}
\left(a^2-\frac{|\mathbf v\times\mathbf a|^2}{c^2}\right)}.
\]
Circular nonrelativistic motion gives cyclotron radiation; relativistic beaming gives synchrotron radiation.
\src{Tong EM \S 6.4}

\subsection{Electromagnetism in Matter}

\subsubsection*{Polarisation, displacement, and bound charge}
Atomic polarisation:
\[
\mathbf p=\alpha\mathbf E,\qquad
\mathbf P=n\langle\mathbf p\rangle,\qquad
\mathbf P=\epsilon_0\chi_e\mathbf E\quad\text{linear dielectric}.
\]
Uniform electron cloud model:
\[
\alpha=4\pi\epsilon_0a^3.
\]
Bound charge:
\[
\boxed{\rho_{\rm b}=-\nabla\cdot\mathbf P},\qquad
\boxed{\sigma_{\rm b}=\mathbf P\cdot\hat{\mathbf n}}.
\]
Displacement:
\[
\boxed{\mathbf D=\epsilon_0\mathbf E+\mathbf P},\qquad
\boxed{\nabla\cdot\mathbf D=\rho_{\rm f}}.
\]
Linear medium:
\[
\mathbf D=\epsilon\mathbf E,\qquad
\epsilon=\epsilon_0(1+\chi_e),\qquad
\epsilon_r=\epsilon/\epsilon_0.
\]
Point charge $Q$ in dielectric sphere:
\[
\mathbf D=\frac{Q}{4\pi r^2}\hat{\mathbf r},\qquad
\mathbf E_{\rm in}=\frac{Q}{4\pi\epsilon r^2}\hat{\mathbf r},\qquad
\mathbf E_{\rm out}=\frac{Q}{4\pi\epsilon_0r^2}\hat{\mathbf r}.
\]
\src{Tong EM \S 7.1}

\subsubsection*{Magnetisation, H field, and bound currents}
\[
\mathbf M=n\langle\mathbf m\rangle,\qquad
\boxed{\mathbf J_{\rm b}=\nabla\times\mathbf M},\qquad
\boxed{\mathbf K_{\rm b}=\mathbf M\times\hat{\mathbf n}}.
\]
\[
\boxed{\mathbf H=\frac{1}{\mu_0}\mathbf B-\mathbf M},\qquad
\boxed{\nabla\times\mathbf H=\mathbf J_{\rm f}}.
\]
Linear magnetic medium:
\[
\mathbf M=\chi_m\mathbf H,\qquad
\mathbf B=\mu\mathbf H,\qquad
\mu=\mu_0(1+\chi_m).
\]
Diamagnet: $-1<\chi_m<0$; paramagnet: $\chi_m>0$; ferromagnet: $\mathbf M\neq0$ at $\mathbf B=0$.
\src{Tong EM \S 7.2}

\subsubsection*{Macroscopic Maxwell equations and boundary data}
Time-dependent bound current:
\[
\boxed{\mathbf J_{\rm b}=\nabla\times\mathbf M+\frac{\pd\mathbf P}{\pd t}}.
\]
\[
\boxed{
\begin{aligned}
\nabla\cdot\mathbf D&=\rho_{\rm f},&
\nabla\times\mathbf H-\frac{\pd\mathbf D}{\pd t}&=\mathbf J_{\rm f},\\
\nabla\cdot\mathbf B&=0,&
\nabla\times\mathbf E+\frac{\pd\mathbf B}{\pd t}&=0.
\end{aligned}}
\]
For $\mathbf D=\epsilon\mathbf E$, $\mathbf B=\mu\mathbf H$, no free sources:
\[
\nabla^2\mathbf E-\frac{1}{v^2}\pd_t^2\mathbf E=0,\qquad
v=\frac{1}{\sqrt{\epsilon\mu}},\qquad n=\frac{c}{v}\simeq\sqrt{\epsilon_r}\quad(\mu\simeq\mu_0).
\]
Plane wave:
\[
\omega^2=v^2k^2,\qquad \mathbf k\cdot\mathbf E_0=\mathbf k\cdot\mathbf B_0=0,\qquad
\mathbf B_0=\frac{\hat{\mathbf k}\times\mathbf E_0}{v}.
\]
Interface normal $\hat{\mathbf n}:1\to2$:
\[
\hat{\mathbf n}\cdot(\mathbf D_2-\mathbf D_1)=\sigma_{\rm f},\qquad
\hat{\mathbf n}\cdot(\mathbf B_2-\mathbf B_1)=0,
\]
\[
\hat{\mathbf n}\times(\mathbf E_2-\mathbf E_1)=0,\qquad
\hat{\mathbf n}\times(\mathbf H_2-\mathbf H_1)=\mathbf K_{\rm f}.
\]
\src{Tong EM \S 7.3}

\subsubsection*{Reflection, refraction, Fresnel coefficients}
Phase matching at a planar interface:
\[
\omega_I=\omega_R=\omega_T,\qquad
k_I\sin\theta_I=k_R\sin\theta_R=k_T\sin\theta_T.
\]
\[
\boxed{\theta_R=\theta_I},\qquad
\boxed{n_1\sin\theta_I=n_2\sin\theta_T}.
\]
For $\mu_1\simeq\mu_2\simeq\mu_0$, $s$ polarisation:
\[
\frac{E_R}{E_I}=\frac{n_1\cos\theta_I-n_2\cos\theta_T}
{n_1\cos\theta_I+n_2\cos\theta_T},\qquad
\frac{E_T}{E_I}=\frac{2n_1\cos\theta_I}
{n_1\cos\theta_I+n_2\cos\theta_T}.
\]
$p$ polarisation:
\[
\frac{E_R}{E_I}=\frac{n_1\cos\theta_T-n_2\cos\theta_I}
{n_1\cos\theta_T+n_2\cos\theta_I},\qquad
\frac{E_T}{E_I}=\frac{2n_1\cos\theta_I}
{n_1\cos\theta_T+n_2\cos\theta_I}.
\]
Brewster and critical angles:
\[
\boxed{\tan\theta_B=\frac{n_2}{n_1}},\qquad
\boxed{\sin\theta_C=\frac{n_2}{n_1}}\quad(n_1>n_2).
\]
For $\theta_I>\theta_C$:
\[
k_{T,x}=\pm i\frac{\omega}{v_2}
\sqrt{\frac{n_1^2}{n_2^2}\sin^2\theta_I-1},
\]
giving an evanescent transmitted field. \src{Tong EM \S 7.4}

\subsubsection*{Dispersion and Kramers-Kronig}
Driven atomic oscillator:
\[
m\ddot{\mathbf r}+m\gamma\dot{\mathbf r}+m\omega_0^2\mathbf r=-q\mathbf E_0e^{-i\omega t},
\qquad
\boxed{\alpha(\omega)=\frac{q^2/m}{\omega_0^2-\omega^2-i\gamma\omega}}.
\]
\[
\epsilon(\omega)=\epsilon_0+n_{\rm at}\alpha(\omega),\qquad
\mathbf D=\epsilon(\omega)\mathbf E.
\]
Complex wavevector:
\[
k^2=\mu\epsilon(\omega)\omega^2,\qquad
\epsilon=\epsilon_1+i\epsilon_2,\qquad k=k_1+ik_2,
\]
\[
\mathbf E(z,t)=\mathbf E(\omega)e^{-k_2z}e^{i(k_1z-\omega t)}.
\]
Velocities:
\[
v_p=\frac{\omega}{k_1},\qquad
v_g=\frac{\dd\omega}{\dd k_1},\qquad
\frac{1}{v_g}=\frac{1}{v_p}+\frac{\omega}{c}\frac{\dd n}{\dd\omega}.
\]
Lorentz oscillator real/imaginary parts:
\[
\epsilon_1=\epsilon_0-\frac{n_{\rm at}q^2}{m}
\frac{\omega^2-\omega_0^2}{(\omega^2-\omega_0^2)^2+\gamma^2\omega^2},
\]
\[
\epsilon_2=\frac{n_{\rm at}q^2}{m}
\frac{\gamma\omega}{(\omega^2-\omega_0^2)^2+\gamma^2\omega^2},\qquad
\omega_p^2=\frac{n_{\rm at}q^2}{m\epsilon_0}.
\]
Causal response:
\[
p(t)=\int_{-\infty}^{\infty}\tilde\alpha(t-t')E(t')\,\dd t',\qquad
\tilde\alpha(t<0)=0
\Longrightarrow \alpha(\omega)\ \text{analytic for }\operatorname{Im}\omega>0.
\]
Kramers-Kronig:
\[
\operatorname{Re}\alpha(\omega)=\mathcal P\int_{-\infty}^{\infty}\frac{\dd\omega'}{\pi}
\frac{\operatorname{Im}\alpha(\omega')}{\omega'-\omega},\qquad
\operatorname{Im}\alpha(\omega)=-\mathcal P\int_{-\infty}^{\infty}\frac{\dd\omega'}{\pi}
\frac{\operatorname{Re}\alpha(\omega')}{\omega'-\omega}.
\]
\src{Tong EM \S 7.5}

\subsubsection*{Conductors and plasma}
Drude equation:
\[
m\frac{\dd\mathbf v}{\dd t}=q\mathbf E-\frac{m}{\tau}\mathbf v,\qquad
\mathbf J=nq\mathbf v=\sigma(\omega)\mathbf E,
\]
\[
\boxed{\sigma(\omega)=\frac{\sigma_{\rm DC}}{1-i\omega\tau}},\qquad
\boxed{\sigma_{\rm DC}=\frac{nq^2\tau}{m}}.
\]
Conducting medium:
\[
\epsilon_{\rm eff}(\omega)=\epsilon(\omega)+i\frac{\sigma(\omega)}{\omega},\qquad
\boxed{k^2=\mu\epsilon_{\rm eff}(\omega)\omega^2}.
\]
Low frequency $\omega\tau\ll1$:
\[
\epsilon_{\rm eff}\simeq i\frac{\sigma_{\rm DC}}{\omega},\qquad
k=\sqrt{\frac{\mu\omega\sigma_{\rm DC}}{2}}(1+i),
\]
\[
\boxed{\delta=\sqrt{\frac{2}{\mu\omega\sigma_{\rm DC}}}},\qquad
\mathbf E(z,t)=\mathbf E_0e^{-z/\delta}e^{i(k_1z-\omega t)}.
\]
High frequency:
\[
\epsilon_{\rm eff}\simeq\epsilon_0\left(1-\frac{\omega_p^2}{\omega^2}\right),\qquad
\omega_p^2=\frac{nq^2}{m\epsilon_0}.
\]
\[
\omega>\omega_p:\ \omega^2=\omega_p^2+c^2k^2,\qquad
\omega<\omega_p:\ k\ \text{imaginary, reflection}.
\]
Longitudinal plasma oscillations:
\[
\epsilon_{\rm eff}(\omega)=0,\qquad
\omega=\omega_p,\qquad \mathbf E\parallel\mathbf k,\qquad \mathbf B=0.
\]
\src{Tong EM \S 7.6}

\subsubsection*{Screening, dielectric function, Thomas-Fermi, Lindhard}
Classical Debye-Huckel:
\[
\nabla^2\phi=-\frac{1}{\epsilon_0}\left[Q\delta^{(3)}(\mathbf r)-\rho_0+\rho(\mathbf r)\right],
\qquad
\rho(\mathbf r)=\rho_0e^{-q\phi/k_BT}.
\]
Linearised:
\[
\left(\nabla^2-\lambda_D^{-2}\right)\phi=-\frac{Q}{\epsilon_0}\delta^{(3)}(\mathbf r),\qquad
\boxed{\lambda_D^2=\frac{\epsilon_0k_BT}{q^2n_0}},
\]
\[
\boxed{\phi(r)=\frac{Q}{4\pi\epsilon_0r}e^{-r/\lambda_D}}.
\]
Static dielectric function:
\[
\phi_{\rm ext}(\mathbf k)=\epsilon(\mathbf k)\phi(\mathbf k),\qquad
\boxed{\epsilon(k)=1-\frac{\rho_{\rm ind}(k)}{\epsilon_0k^2\phi(k)}}.
\]
Debye-Huckel gives
\[
\epsilon(k)=1+\frac{k_D^2}{k^2},\qquad k_D=\lambda_D^{-1}.
\]
Thomas-Fermi:
\[
\rho_{\rm ind}\simeq-q\phi\,\frac{\pd\rho}{\pd\mu},\qquad
\epsilon(k)=1+\frac{q}{\epsilon_0k^2}\frac{\pd\rho}{\pd\mu}.
\]
At $T=0$:
\[
\boxed{\epsilon(k)=1+\frac{q^2g(E_F)}{\epsilon_0k^2}}
\equiv1+\frac{k_{TF}^2}{k^2}.
\]
For 3D nonrelativistic fermions:
\[
g(E)=\frac{g_s}{4\pi^2}\left(\frac{2m}{\hbar^2}\right)^{3/2}E^{1/2},\qquad
k_{TF}^2=\frac{3q\rho_0}{2\epsilon_0E_F}.
\]
Lindhard:
\[
\frac{\rho_{\rm ind}(\mathbf k)}{\phi(\mathbf k)}
=q^2g_s\int\frac{\dd^3k'}{(2\pi)^3}
\frac{f(|\mathbf k+\mathbf k'|)-f(k')}
{E(|\mathbf k+\mathbf k'|)-E(k')},
\]
\[
\boxed{\epsilon(k)=1+\frac{k_{TF}^2}{k^2}F\!\left(\frac{k}{2k_F}\right)},\qquad
\boxed{F(x)=\frac12+\frac{1-x^2}{4x}\log\left|\frac{x+1}{x-1}\right|}.
\]
Nonanalyticity at $k=2k_F$ gives Friedel tail for a screened point charge:
\[
\boxed{\rho_{\rm ind}(r)\approx
\frac{Qk_{TF}^2}{8\pi k_F^2}\frac{\cos(2k_Fr)}{r^3}}.
\]
\src{Tong EM \S 7.7}
